\section*{Limitations}
\label{sec:limitations}

EvLink has several limitations. First, our evaluation focuses on English
Wikipedia-style QA, so its behavior on multilingual, domain-specific, and
rapidly changing corpora remains untested. Second, evidence links depend on
OpenIE quality. Missed, noisy, or underspecified facts can prevent useful
passage transitions or introduce weak links, especially in documents with
implicit relations. Third, EvLink assumes passage-level retrieval units and a
fixed top-five reader context, which may not match long-document summarization
or settings where sentence-level evidence is required. Finally, source-grounded links improve auditability but do not guarantee factual correctness, since the reader can still misinterpret incomplete evidence.

\section*{Ethical Considerations}

This work studies retrieval methods on public QA benchmarks and does
not introduce new human-subject data, private user logs, or personally
identifying annotations. The main ethical risks come from the underlying
corpora and language models: retrieved passages may contain outdated,
biased, or factually incorrect information, and the reader model may
still generate unsupported answers when the evidence is incomplete.
EvLink makes cross-document transitions more auditable by grounding
them in source text, but this should not be interpreted as a guarantee
of factual correctness. Systems based on this method should therefore
include evidence inspection and domain-specific validation before use in
high-stakes settings such as medical, legal, or financial decision
support. We also report auxiliary token costs to make the additional
compute requirements of evidence-link construction transparent.
